% Раздел: Подготовка данных (глава «Метод и экспериментальная установка»)

\subsection{Подготовка данных}
\label{sec:preprocessing}

\textbf{Стилевые изображения.} Восемь эталонных работ организованы в файловую иерархию \texttt{data/styles/\{van\_gogh,monet\}/img\{1..4\}/}, совместимую с параметрами \texttt{folder\_path} и \texttt{caption\_ext} фреймворка \texttt{ai-toolkit}: каждое изображение сопровождается одноимённым \texttt{.txt}-файлом с подписью. Использование четырёх изображений на художника соответствует постановке оригинального B-LoRA~\cite{frenkel2024blora} (для файнтюна достаточно одного референса) и позволяет оценить вариативность результатов внутри стиля при сопоставимом вычислительном бюджете.

Подписи составлены вручную и описывают \textit{исключительно} семантическое содержание~--- предметы, пространственные отношения, освещение~--- без упоминаний имени художника или принадлежности к стилю. Это требование принципиально для B-LoRA: стилевая информация должна кодироваться в весах $\Theta_\text{style}$, а не передаваться через текстовое условие. Например, подпись к «Звёздной ночи»~--- \textit{<<a night sky with swirling clouds and stars over a village with a church steeple>>}; к «Красным виноградникам»~--- \textit{<<a vineyard at sunset with red fields and workers in a rural landscape>>}. В конфигурации обучения задан \texttt{caption\_dropout\_rate: 0.05}: на 5\,\% шагов обучение проводится без текстового условия~\cite{ho2022cfg}, что снижает зависимость адаптера от конкретных формулировок и повышает обобщающую способность.

\textbf{Промпты для оценки.} Сформирован набор из 100~текстовых промптов, извлечённых из \texttt{captions\_val2017.json}: \texttt{seed=42}, длина 5--20~слов, без дубликатов, сохранение в \texttt{data/coco\_prompts.txt}. Размер выбран как компромисс между вычислительной стоимостью генерации (100~прогонов FLUX.1 на эксперимент при 28~шагах инференса) и статистической надёжностью агрегированных метрик; требования~--- семантическое разнообразие сцен и нейтральность по отношению к художественному стилю~\cite{lin2014coco}.

\textbf{Ограничение FID.} Надёжная оценка FID требует тысяч изображений~\cite{heusel2017fid,parmar2022cleanfid}; в работе используется 100~сгенерированных против 500~реальных из ArtBench-10. Вычисляемое значение носит ориентировочный характер и используется для сравнения методов в рамках одного протокола, а не как абсолютная оценка качества.
